% Section A.1, from lectures/L01/appendix/solutions.

\section{Chapter 1: Modern Embedded C++ Concepts}
\label{sol:sec:c1}

% ------------------------------------------------------------------------------------------
\subsection{Exercise Set 1: Default Arguments and Functions}
\label{sol:set:c1.1}
Exercises~\ref{c1:ex:1.1} and~\ref{c1:ex:1.2}, in \file{lectures/L01/appendix/solutions/exercise1}.
Both functions are \code{noexcept}, and both have a default argument, so
\code{debug::log("System started")} logs at level 0 and \code{app::delay_ms()} would wait one
millisecond's worth of loop.

\cppfile{lectures/L01/appendix/solutions/exercise1/main.cpp}

\noindent Output:

\begin{console}
System started, log level = 0
Sensor failure, log level = 2
\end{console}

% ------------------------------------------------------------------------------------------
\subsection{Exercise Set 2: Struct Driver}
\label{sol:set:c1.2}
Exercise~\ref{c1:ex:2.1}, in \file{lectures/L01/appendix/solutions/exercise2}. The driver is a
header, \file{driver/timer.hpp}, included by \file{main.cpp}.

\cppfile{lectures/L01/appendix/solutions/exercise2/driver/timer.hpp}
\cppfile{lectures/L01/appendix/solutions/exercise2/main.cpp}

\noindent Output:

\begin{console}
Creating timer!
Starting timer!
Timeout after 1000 ms!
Timeout after 1000 ms!
Timeout after 1000 ms!
Timeout after 1000 ms!
Timeout after 1000 ms!
Stopping timer!
Destroying timer!
\end{console}

% ------------------------------------------------------------------------------------------
\subsection{Exercise Set 3: References}
\label{sol:set:c1.3}
Exercise~\ref{c1:ex:3.1}, in \file{lectures/L01/appendix/solutions/exercise3}.

\cppfile{lectures/L01/appendix/solutions/exercise3/main.cpp}

\noindent Output:

\begin{console}
Before swap: a = 3, b = 10
After swap: a = 10, b = 3
\end{console}

% ------------------------------------------------------------------------------------------
\subsection{Exercise Set 4: Bit Manipulation Templates}
\label{sol:set:c1.4}
Exercises~\ref{c1:ex:4.1} and~\ref{c1:ex:4.2}, in \file{lectures/L01/appendix/solutions/exercise4}.
Both templates reject a non-integral register type at compile time, and \code{toggle()} iterates
over the parameter pack through a braced initializer list, as \secref{c1:sec:templates} does for
\code{set()}.

\cppfile{lectures/L01/appendix/solutions/exercise4/main.cpp}

\noindent Output:

\begin{console}
Register content: 10101110
\end{console}
